\section{Conclusions}

The segmentation and visualization of the medical data is a basic tool needed for medical professionals to work. The data sensed in this area is full of different body parts but normally to analyse a condition is needed only to visualize a reduced number of organs. This is a complex work that needs to be simplified to properly do the diagnosis and treatment work of the medical professionals. 

To achieve that goal, the program 3D slicer provides different tools to work with the medical data from different points of interest. In this work it has been seen different segmentation tools that the program offers, it has created different measurements that can help the medical professional work and has been created final images for a professional report.

Moreover has been tested different plugins to segmentate the chosen model based on processment technics as the artificial intelligence processment to simplify the segmentation work for different models. It has been proven the existence of a strong community of 3D slicer users that create different plugins to work with different models and objectives and share it worldwide. 

The visualization of the medical data is a basic tool for the medical professionals, with the newest data acquisition systems created it is possible to accurately describe the human body of the patient but more important is to find the adequate information in all the data that will give the key point to know exactly the health state of the patient. That is why simpler and effective tools of segmentation are a central interest point in the medical application. 